\worksheettitle{Наибольшее и наименьшее число}{\modulemark{M06} Версия учителя \textbullet\ ответы и диагностика}

\begin{teacherbox}[Паспорт модуля]
  \textbf{Результат:} ученик составляет максимум и минимум из заданных цифр,
  использует каждую цифру ровно один раз и объясняет выбор через первый
  различающийся разряд. \textbf{Опора:} M05.
\end{teacherbox}

\begin{teacherbox}[Структура ученического маршрута]
  Обязательное ядро включает вход, примеры максимума и минимума, алгоритм,
  задания 1--7, разбор ошибок и контроль. Задания 8--10 выделены как
  дополнительная практика. Самооценка и решение о дальнейшем маршруте в
  ученическом листе оформлены отдельными блоками.
\end{teacherbox}

\begin{checkpointbox}[Вход]
  Полный список: \(307,\ 370,\ 703,\ 730\). Запись \(037\) обозначает число
  37 и не является записью трёхзначного числа. Если ученик ставит ноль первым
  или теряет варианты, сначала верните его к соответствующему заданию M05.
\end{checkpointbox}

\section{Ответы к открытию способа}

Максимум из \(2,\ 5,\ 7,\ 9\): \(9752\). Минимум без нуля: \(2579\).
При одинаковом количестве цифр больше то число, у которого в первом
различающемся разряде стоит большая цифра.

\begin{teacherbox}[Задания 1 и 2]
  Максимумы: \(641;\ 9832;\ 7540;\ 98300;\ 9551\).

  Для минимума \(2579\) первая цифра 2, потому что это наименьшая из данных
  цифр; любой вариант с 5, 7 или 9 в разряде тысяч будет больше.
\end{teacherbox}

\begin{teacherbox}[Минимум с нулём]
  Не используйте неточное правило «поменять ноль со следующей цифрой»: оно
  плохо переносится на несколько нулей. Точный способ: наименьшая ненулевая
  цифра, затем все нули, затем остальные цифры по возрастанию.
\end{teacherbox}

\newpage
\begin{teacherbox}[Задание 3]
  \begin{center}
  \renewcommand{\arraystretch}{1.35}
  \begin{tabular}{c|c|c}
    цифры & максимум & минимум \\
    \hline
    \(1,3,6,8\) & \(8631\) & \(1368\) \\
    \(0,3,6,8\) & \(8630\) & \(3068\) \\
    \(0,1,4,9\) & \(9410\) & \(1049\) \\
    \(0,0,5,7\) & \(7500\) & \(5007\) \\
    \(0,2,2,6\) & \(6220\) & \(2026\) \\
    \(0,0,0,5\) & \(5000\) & \(5000\)
  \end{tabular}
  \end{center}
\end{teacherbox}

\begin{warningbox}[Ловушка с \(3058\)]
  Ученик прав: \(3058\) \textemdash{} минимум. В разряде тысяч цифры равны, а в разряде
  сотен \(0<5\), поэтому \(3058<3508\). Задача проверяет готовность ученика
  обосновать ответ, даже если проверяющий ошибся.
\end{warningbox}

\newpage
\section{Ответы к закреплению}

\begin{teacherbox}[Задание 4]
  \begin{itemize}
    \item максимум \(743\), минимум \(347\);
    \item максимум \(9551\), минимум \(1559\);
    \item максимум \(9620\), минимум \(2069\);
    \item максимум \(98400\), минимум \(40089\);
    \item максимум \(76110\), минимум \(10167\);
    \item максимум \(987432\), минимум \(234789\).
  \end{itemize}
\end{teacherbox}

\begin{teacherbox}[Задание 5]
  \begin{enumerate}
    \item \(841\) и \(814\) различаются в разряде десятков; \(841>814\),
      потому что \(4>1\).
    \item \(207\) и \(270\) различаются в разряде десятков; \(207<270\),
      потому что \(0<7\).
    \item \(40059\) и \(40509\) различаются в разряде сотен;
      \(40059<40509\), потому что \(0<5\).
  \end{enumerate}
\end{teacherbox}

\begin{teacherbox}[Задание 6]
  В числе \(841\) цифра 8 занимает разряд сотен; ни один другой допустимый
  вариант не может иметь в этом разряде большую цифру. При добавлении нуля к
  цифрам \(2,5,9\) максимум становится равен \(9520\): цифры идут по убыванию,
  поэтому ноль занимает последнюю позицию.
\end{teacherbox}

\section{Ошибки и обратные задачи}

\begin{teacherbox}[Задание 7]
  \begin{enumerate}
    \item Лишняя цифра 2; правильный максимум \(742\).
    \item Ноль поставлен первым; правильный минимум \(1069\).
    \item Нарушен порядок после первой ненулевой цифры; минимум \(3058\).
    \item Нарушен порядок по убыванию; максимум \(9721\).
  \end{enumerate}
\end{teacherbox}

\begin{teacherbox}[Задания 8 и 9]
  Для \(8531\): набор \(8,5,3,1\), минимум \(1358\). Для \(9700\): набор
  \(9,7,0,0\), минимум \(7009\).

  В задании 9 верны утверждения 1 и 4. Утверждение 2 требует оговорки о нуле.
  Утверждение 3 неверно: ноль можно использовать, но нельзя ставить первым в
  записи многозначного числа. В задании 10 дан набор \(6,3,0,0\): в максимуме
  нули стоят справа, а в минимуме \textemdash{} после первой ненулевой цифры 3.
\end{teacherbox}

\newpage
\section{Контрольная точка и развилка}

\begin{teacherbox}[Ответы]
  \begin{enumerate}
    \item \(8631\) и \(1368\).
    \item \(8520\) и \(2058\).
    \item \(94400\) и \(40049\).
    \item Правильный максимум \(7520\). Числа \(7520\) и \(7502\) впервые
      различаются в разряде десятков: \(2>0\).
    \item Запись с первым нулём имеет меньше цифр. После выбора наименьшей
      ненулевой первой цифры ноль помещают в следующий справа разряд, потому
      что это даёт наименьшее возможное продолжение записи.
  \end{enumerate}
\end{teacherbox}

\begin{resultbox}[Решение по результату]
  \begin{itemize}
    \item \textbf{Освоено:} все три пары чисел верны, цифры не потеряны,
      объяснение связано с первым различающимся разрядом. Назначьте M06-H.
    \item \textbf{Точечная доработка или M06-H:} одна вычислительная
      перестановка при верном объяснении либо неустойчивое разрядное
      объяснение. Дайте один аналогичный набор цифр или домашнее закрепление.
    \item \textbf{M06-R:} ведущий ноль, потеря цифры, случайный порядок или
      отсутствие разрядного объяснения.
    \item \textbf{Переход к M07 и M06\(\star\):} базовый результат устойчив;
      усложнение выполняется дополнительно.
  \end{itemize}
\end{resultbox}

\begin{teacherbox}[Диагностические различия]
  \begin{itemize}
    \item лишняя или потерянная цифра \textemdash{} не удержано условие «каждая ровно
      один раз»;
    \item ноль первым \textemdash{} не различаются запись числа и количество цифр;
    \item ноль слишком далеко справа \textemdash{} нет поразрядного сравнения;
    \item почти убывающий максимум \textemdash{} порядок выбран «на глаз»;
    \item верный ответ без объяснения \textemdash{} способ ещё не стал осознанным.
  \end{itemize}
\end{teacherbox}

\newpage
\section{Ответы к дополнительным листам}

\begin{teacherbox}[M06-R]
  Для \(0,3,4,9\): максимум \(9430\), минимум \(3049\). В сравнении
  \(3049<3409\) первый различающийся разряд \textemdash{} сотни; \(0<4\).

  Повторная проверка на новом наборе \(0,3,3,7\): максимум \(7330\), минимум
  \(3037\). На второй странице для \(0,0,5,7\): максимум \(7500\), минимум
  \(5007\); для \(0,2,2,6\): максимум \(6220\), минимум \(2026\).
\end{teacherbox}

\begin{teacherbox}[M06\(\star\)]
  Пример: второе по величине число после \(9752\) \textemdash{} \(9725\).

  Таблица: максимум \(631\), второе число \(613\); максимум \(8520\), второе
  число \(8502\); максимум \(9441\), второе число \(9414\); максимум \(7300\),
  второе число \(7030\).
  Для \(1,4,4,9\) перестановка одинаковых четвёрок не меняет число; первое
  действительное уменьшение даёт \(9414\).

  Дополнительные задания: максимум \(76430\), второе число \(76403\);
  максимум \(99520\), второе число \(99502\). При доказательстве ученик должен
  сохранить максимально возможные одинаковые начальные цифры и показать, что
  меньшего изменения в оставшейся части нет.
\end{teacherbox}

\begin{notebox}[Домашняя работа]
  M06-H назначается после контрольной точки. При проверке оценивайте отдельно:
  верность числа, использование всех цифр и разрядное объяснение.
\end{notebox}
